
\subsubsection{2.1 Consistency-Based Uncertainty Quantification}

Self-consistency as a reliability signal was established by Wang et al. (2022), who showed that sampling N outputs and selecting the most frequent answer reliably improves reasoning task performance --- demonstrating that output diversity correlates with uncertainty. Manakul et al. (2023) operationalized this as SelfCheckGPT: a black-box hallucination detection method that scores consistency across N samples using NLI-based contradiction detection, BERTScore, or n-gram overlap. The NLI variant achieves NonFact AUC-PR of 92.50 vs. 83.21 for log-probability signals on WikiBio, establishing that consistency-based signals outperform log-prob signals in document generation tasks.

The foundational formalization of consistency-based UQ is Semantic Entropy (SE) \citep{Kuhn2023Semantic}, which defines uncertainty as the entropy over semantic equivalence classes: H(p(C|x)) = $-\Sigma$\_C p(C|x) log p(C|x), where p(C|x) = $\Sigma\_{s\in C}$ p(s|x) sums token probabilities of all outputs in the same NLI-defined cluster. Kuhn et al. demonstrate that SE outperforms raw token entropy on TriviaQA and BioASQ (AUROC 0.83 on TriviaQA with OPT-30B). The present work directly builds on this formulation (N=5, bidirectional DeBERTa-MNLI NLI clustering).

Nguyen et al. (2025) (SNNE) introduced pairwise SE computation to handle multi-sentence answers where the original SE formulation weakens. This directly motivates the present work's scope restriction to short-answer factual QA (TriviaQA, $\leq 50$ token answers), where original SE remains the strongest consistency-based signal. Bouchard and Chauhan (2025) (UQLM) benchmark 50+ UQ methods across 24 model-dataset combinations and find that NLI-based black-box scorers are best in 13/24 AUROC scenarios, providing strong external validation that SE-family signals are competitive.

\textbf{Gap:} None of these works directly measure the Pearson correlation between SE and min\_logprob on open-weight LLMs. They report standalone AUROC comparisons, not cross-signal statistical independence. The present work closes this gap.

\subsubsection{2.2 Token Log-Probability Signals for Hallucination Detection}

Token log-probabilities are the most computationally efficient UQ signal class, requiring only a single forward pass. Malinin and Gales (2021) formalize sequence-level uncertainty measures from log-probabilities. Min\_logprob has been empirically shown to achieve strong single-pass baseline performance on TriviaQA dev with Llama-3.1-8B (AUROC approximately 0.825; internal pipeline experiment h-m1).

The relationship between first-token confidence and SE was characterized by Gabriel (2026), who finds r = 0.54--0.76 between first-token probability and semantic agreement in Llama-3.1-8B, TriviaQA --- suggesting that for first-token signals, moderate correlation with SE exists. A critical distinction: min\_logprob is not first-token confidence. It is the minimum over all tokens in greedy decode, capturing arbitrarily late positions. This operational difference drives near-zero Pearson r in the measurements reported here (|r| = 0.026).

Dey et al. (2025) (UAF Ensemble) show that fusing UQ signals improves factual accuracy by 8\%, further motivating the ensemble direction.

\textbf{Gap:} No prior work directly measures the statistical independence of min\_logprob and SE on open-weight LLMs under a controlled evaluation protocol. The |r|=0.026 measurement reported here fills this gap.

\subsubsection{2.3 LM-as-a-Judge Evaluation Methodology}

Santilli et al. (2025) show that AUROC rankings for hallucination detection methods are systematically distorted by response length bias --- ROUGE-L evaluation is most severely affected. They recommend LM-as-a-judge with cross-model design as the most human-aligned correctness function. Xiong et al. (2023) benchmark consistency-based signals as the most reliable black-box UQ method overall.

The present work implements the cross-model judge recommendation directly: Qwen-2.5-7B-Instruct evaluates Llama-3.1-8B-Instruct outputs on TriviaQA dev. The resulting Spearman $\rho$(SE, judge) = $-0.026$ directly validates this design choice, satisfying the circularity threshold of |$\rho$| < 0.40 with substantial margin.

